% Pursuing Stacks (cotes 134-2 to 134-8): archive cote 134-6, site page 4 (PDF page 5), typescript page 323.
% typescript. Draft compared: ps5.tex:1521-1588.
% AI-transcribed by Claude Opus 5.5 (claude-opus-5-5, adaptive thinking, effort high), Message Batches API, batch ps3 (msgbatch_01XdFmuVVBAcgc7EFxdtCqsj), 2026-11-12; not yet checked by eye.
\scanpage{134-6/4}
\mspage{323}

\xout{sectionx82x(px272)} the corresponding functor
\[ \hand{\overline{i^*}} : (\mathrm{Hot})_{\hand{\deleted{\illegible{1}}}} \hand{\longrightarrow} \mathrm{H}\unsure{\over{e}{o}|o}\mathrm{t}_{\hand{\Delta}} \]
is independant of such choice\inserted{,} up to canonical isomorphism.

Instead of \hand{$\Delta$} , we could have worked with \hand{$\Box$} \unsure{of|or} \hand{$\Diamond$} instead, as
these give rise to a Dold-Puppe functor (which is still an equivalence),
and (almost certainly, see below) to a corresponding variant of the
``Whitehead theorem''. We thus get two other ways for defining a canoni-%
cal ``abelianization functor'' (1) for homotopy types, and it should be
an easy and pleasant exercise to show these three functors are canonical%
ly \xout{equi} isomorphic, using the fact that a product of test categories
is again a test category.

\inserted{\uline{Remark.}} Of course, when concerned merely with defining a functor (1)\unsure{;}
we do\unsure{n't|'nt} really need the Whitehead and Dold-Puppe theorems - indeed
instead of taking the functor
\[ \mathrm{DP} \circ \mathrm{Wh}_{\hand{\Delta}} : \hand{\Delta}^{\wedge} \longrightarrow \mathrm{Ch.(Ab)} \quad , \]
we could have taken directly (using the standard boundary operations
between the components of a semisimplicial abelian group) the standard
\xout{semi} chain complex structure of $\underline{Z}^{(K.)}$ (for K. in \hand{$\Delta$}$^{\wedge}$) , without
taking the trouble and normalizing it à la Dold-Puppe -- and it is
a lot more trivial than either Whitehead's or Dold-Puppe's theorem,
that the latter functor transforms \xout{quasi} weak equivalences into
quasi-isomorphisms ; \xout{and} moreover it gives rise to the same functor
\[ (5) \qquad \mathrm{Hot}_{\hand{\Delta}} \longrightarrow \mathrm{D.(Ab)}\hand{(=\mathrm{Hotab})} \]
as $\mathrm{DP}\circ\mathrm{Wh}_{\hand{\Delta}}$ .

Let now A be \uline{any} small category, we are interested in the
functor
\[ (6) \qquad \mathrm{Hot}_{A} \longrightarrow \mathrm{D.(Ab)}\hand{(=\mathrm{Hotab})} \]
obtained as the composition
\[ (7) \qquad \mathrm{Hot}_{A} \longrightarrow (\mathrm{Hot}) \longrightarrow \mathrm{D.(Ab)} \quad , \]
where the second functor is the abelianization functor (1), and the
f\over{o}{i}rst is the canonical functor, deduced by localization from
\[ (8) \qquad i_{A} : A^{\wedge} \longrightarrow (\mathrm{Cat}) \quad , \quad a \longmapsto A_{/a} \ . \]
We see immediately that for $A = $ \hand{$\Delta$} , the functor (6) reduces to
(5) up to canonical isomorphism - and the same of course \xout{for} when

\margin{4}
